% !TEX program = lualatex
\documentclass[a4paper,10pt,twocolumn]{article}
\usepackage{kaitoushu}
\renewcommand{\baselinestretch}{1.2}
\lhead{応用数学 問題集}
\problemrange{問93〜103, 118〜124}

\begin{document}
\thispagestyle{fancy}

\begin{center}
  {\Large \textbf{応用数学　2章　Check　解答}}
\end{center}
\vspace{0.5em}

\noindent\textbf{\large【Check】}
\vspace{0.3em}

\mondai{93}{
次の関数のラプラス変換を求めよ．
(1) $4t^2-3t+2$ \quad (2) $(t-1)^3$ \quad (3) $4e^{3t}-3e^{2t}$ \quad (4) $e^{3t+2}$
}

\kotae{
(1) $\mathcal{L}[4t^2-3t+2] = \dfrac{8}{s^3}-\dfrac{3}{s^2}+\dfrac{2}{s} = \dfrac{8-3s+2s^2}{s^3} = \dfrac{2s^2-3s+8}{s^3} \quad (s>0)$

(2) $(t-1)^3 = t^3-3t^2+3t-1$ より $\mathcal{L} = \dfrac{6}{s^4}-\dfrac{6}{s^3}+\dfrac{3}{s^2}-\dfrac{1}{s} = \dfrac{6-6s+3s^2-s^3}{s^4} = \dfrac{-s^3+3s^2-6s+6}{s^4} \quad (s>0)$

(3) $\mathcal{L}[4e^{3t}-3e^{2t}] = \dfrac{4}{s-3}-\dfrac{3}{s-2} = \dfrac{4(s-2)-3(s-3)}{(s-3)(s-2)} = \dfrac{s+1}{(s-3)(s-2)} \quad (s>3)$

(4) $e^{3t+2}=e^2 e^{3t}$ より $\mathcal{L} = \dfrac{e^2}{s-3} \quad (s>3)$
}

\mondai{94}{
定義に従って，関数 $te^{\alpha t}$ のラプラス変換を求めよ．ただし，$\alpha$ は定数とする．
}

\kotae{
$\mathcal{L}[te^{\alpha t}] = \displaystyle\int_0^\infty te^{\alpha t}e^{-st}\,dt = \int_0^\infty te^{-(s-\alpha)t}\,dt$

部分積分より $= \dfrac{1}{(s-\alpha)^2} \quad (s>\alpha)$
}

\mondai{95}{
$\mathcal{L}[\sinh \omega t]$ を求めよ．ただし，$\omega$ は $0$ でない定数とする．
}

\kotae{
$\mathcal{L}[\sinh \omega t] = \dfrac{1}{2}\!\left(\dfrac{1}{s-\omega}-\dfrac{1}{s+\omega}\right) = \dfrac{\omega}{s^2-\omega^2} \quad (s>|\omega|)$
}

\mondai{96}{
次の関数 $f(t)$ を単位ステップ関数を用いて表せ．また，$\mathcal{L}[f(t)]$ を求めよ．
(1) $f(t) = \begin{cases} 1 & (0 < t < 2,\ t \geq 3) \\ 0 & (2 \leq t < 3) \end{cases}$
(2) $f(t) = \begin{cases} 1 & (0 < t < 1) \\ 0 & (1 \leq t < 3) \\ 2 & (t \geq 3) \end{cases}$
}

\kotae{
(1) $f(t) = 1 - U(t-2) + U(t-3)$

$\mathcal{L}[f(t)] = \dfrac{1}{s} - \dfrac{e^{-2s}}{s} + \dfrac{e^{-3s}}{s} = \dfrac{1 - e^{-2s} + e^{-3s}}{s} \quad (s>0)$

(2) $f(t) = 1 - U(t-1) + 2U(t-3)$

$\mathcal{L}[f(t)] = \dfrac{1}{s} - \dfrac{e^{-s}}{s} + \dfrac{2e^{-3s}}{s} = \dfrac{1 - e^{-s} + 2e^{-3s}}{s} \quad (s>0)$
}

\mondai{97}{
次の関数のラプラス変換を求めよ．
(1) $te^{3t}$ \quad (2) $t^2e^{-2t}$ \quad (3) $e^{2t}\sin 3t$ \quad (4) $e^{-t}\cos 2t$
}

\kotae{
(1) $\mathcal{L}[te^{3t}] = \dfrac{1}{(s-3)^2} \quad (s>3)$

(2) $\mathcal{L}[t^2 e^{-2t}] = \dfrac{2}{(s+2)^3} \quad (s>-2)$

(3) $\mathcal{L}[e^{2t}\sin 3t] = \dfrac{3}{(s-2)^2+9} \quad (s>2)$

(4) $\mathcal{L}[e^{-t}\cos 2t] = \dfrac{s+1}{(s+1)^2+4} \quad (s>-1)$
}

\mondai{98}{
関数 $f(t) = (t-2)^2 U(t-2)$ のグラフをかけ．また，$\mathcal{L}[f(t)]$ を求めよ．
}

\kotae{
グラフ：$t \geq 2$ で放物線 $(t-2)^2$（$t<2$ で $0$）．

$\mathcal{L}[(t-2)^2 U(t-2)] = \dfrac{2e^{-2s}}{s^3} \quad (s>0)$
}

\mondai{99}{
像関数の微分法則を用いて，$\mathcal{L}[t\sinh \omega t]$，$\mathcal{L}[t\cosh \omega t]$ を求めよ．ただし，$\omega$ は $0$ でない定数とする．
}

\kotae{
$\mathcal{L}[\sinh \omega t] = \dfrac{\omega}{s^2-\omega^2}$ の微分法則より

$\mathcal{L}[t\sinh \omega t] = -\dfrac{d}{ds}\dfrac{\omega}{s^2-\omega^2} = \dfrac{2\omega s}{(s^2-\omega^2)^2} \quad (s>|\omega|)$

$\mathcal{L}[t\cosh \omega t] = -\dfrac{d}{ds}\dfrac{s}{s^2-\omega^2} = \dfrac{s^2+\omega^2}{(s^2-\omega^2)^2} \quad (s>|\omega|)$
}

\mondai{100}{
関数 $f(t)$ が $f'(t) + 3f(t) = \sin t$，$f(0) = 1$ を満たすとき，原関数の微分法則を用いて，$\mathcal{L}[f(t)]$ を求めよ．
}

\kotae{
$f'(t)+3f(t)=\sin t$，$f(0)=1$ をラプラス変換すると
$sF(s)-1+3F(s) = \dfrac{1}{s^2+1}$

$(s+3)F(s) = 1+\dfrac{1}{s^2+1} = \dfrac{s^2+2}{s^2+1}$

$\mathcal{L}[f(t)] = F(s) = \dfrac{s^2+2}{(s+3)(s^2+1)} \quad (s>-3)$
}

\mondai{101}{
像関数の高次微分法則を用いて，$\mathcal{L}[t^2 \cos t]$ を求めよ．
}

\kotae{
$\mathcal{L}[\cos t] = \dfrac{s}{s^2+1}$ の2階微分法則より

$\mathcal{L}[t^2\cos t] = \dfrac{d^2}{ds^2}\dfrac{s}{s^2+1} = \dfrac{2s(s^2-3)}{(s^2+1)^3} \quad (s>0)$
}

\mondai{102}{
次の関数の逆ラプラス変換を求めよ．
(1) $\dfrac{1}{(s-2)^4}$ \quad
(2) $\dfrac{s+2}{s^2-2s+1}$ \quad
(3) $\dfrac{s}{s^2-s-6}$
(4) $\dfrac{2s+3}{s^2+4}$ \quad
(5) $\dfrac{2s+1}{s^2+2s+5}$
}

\kotae{
(1) $\dfrac{1}{(s-2)^4}$ より $\mathcal{L}^{-1} = \dfrac{1}{6}t^3 e^{2t}$

(2) $\dfrac{s+2}{(s-1)^2} = \dfrac{1}{s-1}+\dfrac{3}{(s-1)^2}$

$\mathcal{L}^{-1} = e^t + 3te^t$

(3) $\dfrac{s}{(s-3)(s+2)} = \dfrac{3}{5}\cdot\dfrac{1}{s-3}+\dfrac{2}{5}\cdot\dfrac{1}{s+2}$

$\mathcal{L}^{-1} = \dfrac{3}{5}e^{3t}+\dfrac{2}{5}e^{-2t}$

(4) $\dfrac{2s+3}{s^2+4} = 2\cdot\dfrac{s}{s^2+4}+\dfrac{3}{2}\cdot\dfrac{2}{s^2+4}$ より $\mathcal{L}^{-1} = 2\cos 2t+\dfrac{3}{2}\sin 2t$

(5) $\dfrac{2s+1}{(s+1)^2+4} = \dfrac{2(s+1)-1}{(s+1)^2+4}$

$\mathcal{L}^{-1} = 2e^{-t}\cos 2t - \dfrac{1}{2}e^{-t}\sin 2t$
}

\mondai{103}{
次の関数の逆ラプラス変換を求めよ．
(1) $\dfrac{s-5}{(s+1)(s-2)(s-3)}$ \quad
(2) $\dfrac{2s^2+7}{(s-2)^2(s+3)}$
}

\kotae{
(1) 部分分数分解：$\dfrac{s-5}{(s+1)(s-2)(s-3)} = \dfrac{A}{s+1}+\dfrac{B}{s-2}+\dfrac{C}{s-3}$

$A=\dfrac{-6}{(-3)(-4)}=-\dfrac{1}{2}$，$B=\dfrac{-3}{(3)(-1)}=1$，$C=\dfrac{-2}{(4)(1)}=-\dfrac{1}{2}$

$\mathcal{L}^{-1} = -\dfrac{1}{2}e^{-t}+e^{2t}-\dfrac{1}{2}e^{3t}$

(2) $\dfrac{2s^2+7}{(s-2)^2(s+3)} = \dfrac{A}{s-2}+\dfrac{B}{(s-2)^2}+\dfrac{C}{s+3}$

$B=3$，$C=1$，$A=1$

$\mathcal{L}^{-1} = e^{2t}+3te^{2t}+e^{-3t}$
}

\mondai{118}{
次の微分方程式を解け．
(1) $\dfrac{dx}{dt} - x = te^t \quad (t=0 \text{ のとき } x=1)$
(2) $\dfrac{d^2x}{dt^2} + 4x = t \quad (t=0 \text{ のとき } x=1,\ \dfrac{dx}{dt}=2)$
(3) $\dfrac{d^2x}{dt^2} - 5\dfrac{dx}{dt} + 6x = e^{3t} \quad (t=0 \text{ のとき } x=1,\ \dfrac{dx}{dt}=5)$
}

\kotae{
(1) $(s-1)X = 1+\dfrac{1}{(s-1)^2}$ より $X = \dfrac{1}{s-1}+\dfrac{1}{(s-1)^3}$

$x(t) = e^t + \dfrac{1}{2}t^2 e^t$

(2) $(s^2+4)X = s+2+\dfrac{1}{s^2}$

$X = \dfrac{s}{s^2+4}+\dfrac{2}{s^2+4}+\dfrac{1}{s^2(s^2+4)}$

$\dfrac{1}{s^2(s^2+4)} = \dfrac{1}{4}\!\left(\dfrac{1}{s^2}-\dfrac{1}{s^2+4}\right)$

$x(t) = \cos 2t + \sin 2t + \dfrac{1}{4}t - \dfrac{1}{8}\sin 2t = \cos 2t + \dfrac{7}{8}\sin 2t + \dfrac{t}{4}$

(3) $(s^2-5s+6)X = s+\dfrac{1}{s-3}$，$(s-2)(s-3)X = s+\dfrac{1}{s-3}$

$X = \dfrac{s}{(s-2)(s-3)}+\dfrac{1}{(s-3)^2(s-2)}$

$\dfrac{s}{(s-2)(s-3)} = -\dfrac{2}{s-2}+\dfrac{3}{s-3}$，$\dfrac{1}{(s-3)^2(s-2)} = \dfrac{1}{s-2}-\dfrac{1}{s-3}+\dfrac{1}{(s-3)^2}$

$x(t) = -e^{2t}+2e^{3t}+te^{3t}$
}

\mondai{119}{
微分方程式 $\dfrac{d^2x}{dt^2} + 9x = \cos t$ について，次の問いに答えよ．
(1) 初期条件 $x(0)=1$，$x'(0)=1$ を満たす解を求めよ．
(2) 境界条件 $x(0)=1$，$x\!\left(\dfrac{\pi}{4}\right)=0$ を満たす解を求めよ．
(3) 一般解を求めよ．
}

\kotae{
(1) $(s^2+9)X = s+1+\dfrac{s}{s^2+1}$

$X = \dfrac{s+1}{s^2+9}+\dfrac{s}{(s^2+1)(s^2+9)}$

$\dfrac{s}{(s^2+1)(s^2+9)} = \dfrac{1}{8}\!\left(\dfrac{s}{s^2+1}-\dfrac{s}{s^2+9}\right)$

$x(t) = \cos 3t+\dfrac{1}{3}\sin 3t+\dfrac{1}{8}\cos t-\dfrac{1}{8}\cos 3t$
$= \dfrac{7}{8}\cos 3t+\dfrac{1}{3}\sin 3t+\dfrac{1}{8}\cos t$

(2) 一般解 $x(t) = C_1\cos 3t + C_2\sin 3t + \dfrac{1}{8}\cos t$ に $x(0)=1$ を代入すると $C_1 = \dfrac{7}{8}$

$x\!\left(\dfrac{\pi}{4}\right) = \dfrac{7}{8}\cos\dfrac{3\pi}{4}+C_2\sin\dfrac{3\pi}{4}+\dfrac{1}{8}\cos\dfrac{\pi}{4} = 0$

$-\dfrac{7\sqrt{2}}{16}+\dfrac{\sqrt{2}}{2}C_2+\dfrac{\sqrt{2}}{16} = 0$ より $C_2 = \dfrac{3}{4}$

$x(t) = \dfrac{7}{8}\cos 3t+\dfrac{3}{4}\sin 3t+\dfrac{1}{8}\cos t$

(3) 一般解：$x(t) = C_1\cos 3t + C_2\sin 3t + \dfrac{1}{8}\cos t$
}

\mondai{120}{
微分方程式 $\dfrac{d^2x}{dt^2} + 2\dfrac{dx}{dt} + 5x = 0$ について，次の問いに答えよ．
(1) 初期条件 $x(0)=1$，$x'(0)=2$ を満たす解を求めよ．
(2) 境界条件 $x(0)=1$，$x\!\left(\dfrac{\pi}{4}\right)=0$ を満たす解を求めよ．
(3) 一般解を求めよ．
}

\kotae{
(1) $(s^2+2s+5)X = s+4$ より $X = \dfrac{s+4}{(s+1)^2+4} = \dfrac{(s+1)+3}{(s+1)^2+4}$

$x(t) = e^{-t}\cos 2t + \dfrac{3}{2}e^{-t}\sin 2t$

(2) 一般解 $x(t) = e^{-t}(C_1\cos 2t + C_2\sin 2t)$ に $x(0)=1$ を代入すると $C_1 = 1$

$x\!\left(\dfrac{\pi}{4}\right) = e^{-\pi/4}(1\cdot\cos\dfrac{\pi}{2}+C_2\sin\dfrac{\pi}{2}) = e^{-\pi/4}C_2 = 0$ より $C_2 = 0$

$x(t) = e^{-t}\cos 2t$

(3) 一般解：$x(t) = e^{-t}(C_1\cos 2t + C_2\sin 2t)$
}

\mondai{121}{
次のたたみ込みを求めよ．
(1) $\cos t * \cos t$ \quad (2) $t * e^{2t}$ \quad (3) $t^2 * \sin t$
}

\kotae{
(1) $\cos t * \cos t = \displaystyle\int_0^t \cos\tau\cos(t-\tau)\,d\tau$

積和公式 $\cos\tau\cos(t-\tau)=\dfrac{1}{2}[\cos t+\cos(2\tau-t)]$ を用いて
$= \dfrac{1}{2}t\cos t + \dfrac{1}{2}\sin t$

(2) $t * e^{2t} = \displaystyle\int_0^t \tau e^{2(t-\tau)}\,d\tau = e^{2t}\int_0^t \tau e^{-2\tau}\,d\tau$

$= \dfrac{1}{4}(e^{2t}-2t-1)$

(3) $\mathcal{L}[t^2*\sin t] = \dfrac{2}{s^3}\cdot\dfrac{1}{s^2+1} = \dfrac{2}{s^3(s^2+1)}$

部分分数分解：$\dfrac{2}{s^3(s^2+1)} = -\dfrac{2}{s}+\dfrac{2}{s^3}+\dfrac{2s}{s^2+1}$

$t^2 * \sin t = -2 + t^2 + 2\cos t$
}

\mondai{122}{
次の積分方程式を満たす関数 $x(t)$ を求めよ．
(1) $\int_0^t x(t-\tau)\sin 2\tau\,d\tau = t\sin 2t$
(2) $x(t) - 2\int_0^t x(t-\tau)\cos\tau\,d\tau = \sin t$
(3) $x'(t) + 2x(t) + 4\int_0^t x(t-\tau)e^{2\tau}\,d\tau = t^2$，$x(0)=0$
}

\kotae{
(1) たたみ込みのラプラス変換：$X(s)\cdot\dfrac{2}{s^2+4} = \mathcal{L}[t\sin 2t] = \dfrac{4s}{(s^2+4)^2}$

$X(s) = \dfrac{4s}{(s^2+4)^2}\cdot\dfrac{s^2+4}{2} = \dfrac{2s}{s^2+4}$

$x(t) = 2\cos 2t$

(2) $X(s) - 2X(s)\cdot\dfrac{s}{s^2+1} = \dfrac{1}{s^2+1}$

$X(s)\!\left(1-\dfrac{2s}{s^2+1}\right) = \dfrac{1}{s^2+1}$

$X(s) = \dfrac{1}{(s-1)^2}$ より $x(t) = te^t$

(3) $sX+2X+4X\cdot\dfrac{1}{s-2} = \dfrac{2}{s^3}$

$X\!\left(s+2+\dfrac{4}{s-2}\right) = \dfrac{2}{s^3}$，$X\cdot\dfrac{s^2}{s-2} = \dfrac{2}{s^3}$

$X = \dfrac{2(s-2)}{s^5} = \dfrac{2}{s^4}-\dfrac{4}{s^5}$

$x(t) = \dfrac{t^3}{3}-\dfrac{t^4}{6}$
}

\mondai{123}{
微分方程式 $y'' + 4y = x(t)$，$y(0) = 0$，$y'(0) = 0$ で表される線形システムの伝達関数を求めよ．また，出力 $y(t)$ を入力 $x(t)$ を用いて表せ．
}

\kotae{
$(s^2+4)Y = X$ より $G(s) = \dfrac{1}{s^2+4}$

$g(t) = \dfrac{1}{2}\sin 2t$

$y(t) = \dfrac{1}{2}\displaystyle\int_0^t x(\tau)\sin 2(t-\tau)\,d\tau$
}

\mondai{124}{
次の微分方程式の解を求めよ．
$y'' + y' - 6y = \delta(t)$，$y(0) = 0$，$y'(0) = 0$
}

\kotae{
$(s^2+s-6)Y = 1$ より $Y = \dfrac{1}{(s+3)(s-2)} = \dfrac{1}{5}\!\left(\dfrac{1}{s-2}-\dfrac{1}{s+3}\right)$

$y(t) = \dfrac{1}{5}(e^{2t}-e^{-3t})$
}

\end{document}